\documentclass[11pt]{exam}
\pagestyle{headandfoot}
\firstpageheader{}{}{}
\runningheader{\date}{}{\class, \teacher}
\runningheadrule
\firstpagefooter{}{}{}
\runningfooter{}{Viel Erfolg!}{Seite \thepage\:von \numpages}
\runningfootrule
%=========================================================
\newcommand{\theme}{Geraden \& Ebenen}
\newcommand{\teacher}{Jonas Guler}
\newcommand{\class}{4GaLaWa}
\newcommand{\serie}{Repetitionsprüfung}
%\newcommand{\date}{\today}
\newcommand{\duration}{45 Minuten}
\newcommand{\name}{\bf Name:}
\newcommand{\grade}{Note:}
\newcommand{\datum}{06.05.24}
%=========================================================
\begin{document}
\fontsize{14}{14}\selectfont
%========================================================

%========================================================
\begin{tabular*}{\textwidth}{l @{\extracolsep{\fill}}l @{\extracolsep{6pt}}l}
\textbf{\theme} & \textbf{\grade} & {\answerbox}\\
\textbf{Klasse \class} & & \\
\textbf{\teacher}  & {\name} &\textbf{\hrulefill}\\[8pt]
%-------------------
\end{tabular*}
%--------------------
\begin{center}
\rule[1ex]{\textwidth}{1pt}
Dauer: \duration \hspace{9cm}Datum: \datum\\
\vspace{5pt}
\rule[2ex]{\textwidth}{1pt}
\end{center}
%=========================================================

%=========================================================
\begin{center}
\textbf{Bewertungstabelle}\\
\small Wird von der Lehrperson ausgefüllt\\
\vspace{5pt}
\addpoints
\gradetable[h][questions]\\
\vspace{5pt}
Der Rechenweg muss \textbf{vollständig und nachvollziehbar} sein.\\
Lösen Sie die Prüfung mit einen schwarzen oder blauen Kugelschreiber. Resultate in Bleistift, roter oder grüner Farbe werden nicht berücksichtigt.\\
Als Hilfsmittel sind der grosse Taschenrechner und das Formelbüchlein erlaubt.
\end{center}

\noindent
\rule[1ex]{\textwidth}{1pt}

\vspace{10cm}
\begin{center}
    %\Huge{Serie \serie}
\end{center}

\newpage
\begin{questions}
%=========================================================
%-------------------TESTFRAGEN
%=========================================================
\noaddpoints
\question[\totalpoints]
\addpoints
\begin{parts}
    \part[1] Erkläre einem Laien in etwa zwei Sätzen, wie du entscheidest, ob ein Punkt $P$ auf einer Gerade $g$ liegt.
\end{parts}
\vspace{5pt}
\antwort{4}
\begin{parts}
    \setcounter{partno}{1}
    \part[2] Welche der folgenden Geraden passt nicht dazu? Begründe deine Antwort.
    \[g:\Vec{x}{y}{z}=\Vec{3}{-5}{6}+t\*\Vec{6}{11}{-3},\quad h:\Vec{x}{y}{z}=\Vec{9}{6}{3}+t\*\Vec{-6}{-11}{3}\]
    \[l:\Vec{x}{y}{z}=\Vec{3}{-5}{6}+t\*\Vec{-24}{-44}{12},\quad m:\Vec{x}{y}{z}=\Vec{9}{6}{3}+t\*\Vec{2}{4}{-1}\]
\end{parts}
\vspace{5pt}
\antwort{11.5}

%---------------------------------------------------------
\newpage
\noaddpoints
\question[\totalpoints]
\addpoints
\begin{parts}
    \part[3] Zeige, dass die beiden Geraden $g$ und $h$ parallel sind.
        \[g:\Vec{x}{y}{z}=\Vec{2}{2}{2}+s\*\Vec{9}{-3}{6},\quad h:\Vec{x}{y}{z}=\Vec{12}{-\frac{3}{2}}{9}+t\*\Vec{-6}{2}{-4}\]
\end{parts}
\vspace{5pt}
\antwort{10}
\begin{parts}
    \setcounter{partno}{1}
    \part[2] Gib die Parametergleichung der Ebene $E$ an, die die beiden Geraden $g$ und $h$ enthält.
\end{parts}
\vspace{5pt}
\antwort{8}

%---------------------------------------------------------
\newpage
\question[4]
Berechne den Schnittpunkt und den Schnittwinkel der beiden Geraden $g$ und $h$.
\[g:\Vec{x}{y}{z}=\Vec{1}{-2}{-5}+s\*\Vec{4}{5}{6},\quad h:\Vec{x}{y}{z}=\Vec{-3}{0}{3}+t\*\Vec{3}{2}{1}\]
\vspace{5pt}
\antwort{8}
\question[2]
Entscheide in der Tabelle unten, ob die Aussagen wahr oder falsch sind. Die Antworten müssen nicht begründet werden.\newline
\begin{small}
    Bewertung: Für $0-2$ korrekte Antworten gibt es keine Punkte; $3$ korrekte Antworten einen Punkt und falls alle Antworten korrekt sind, gibt es $2$ Punkte.
\end{small}
\begingroup
\renewcommand{\arraystretch}{1.5}
\begin{table}[h]
    \centering
    \begin{tabular}{|p{12cm}|c|c|}\hline
         & wahr & falsch \\\hline
        Der Punkt $P(1,6,-3)$ liegt auf der Ebene $E:5x-y-2z-5=0$. & & \\\hline
        Falls das Gleichungssystem zur Berechnung eines Durchstosspunktes keine Lösung besitzt, ist die Gerade vollständig in der Ebene enthalten. & & \\\hline
        Zwei windschiefe Geraden besitzen keine gemeinsamen Punkte. & & \\\hline
        \begingroup\renewcommand{\arraystretch}{1} Die Gerade $g:\Vec{x}{y}{z}=\Vec{0}{0}{7}+t\*\Vec{-6}{2}{4}$ enthält den Punkt $P(-21,7,21)$.\endgroup & & \\\hline
    \end{tabular}
\end{table}
\endgroup

%---------------------------------------------------------
\newpage
\question[2]
Gib die Koordinatengleichung der Ebene $F$ an, die parallel zur Ebene $E$ liegt und den Punkt $P$ enthält.
\[E:2x+y-3z-37=0\quad\text{und}\quad P(4,8,7)\]
\vspace{5pt}
\antwort{7}
\vfill
\question[3]
Berechne den Durchstosspunkt der Ebene $E$ mit der Gerade $g$.
\[g:\Vec{x}{y}{z}=\Vec{7}{6}{3}+t\*\Vec{2}{-1}{-2}, \quad E:2x-2y-z+17=0\]
\vspace{5pt}
\antwort{7}

%---------------------------------------------------------
\newpage
\question[4] %Matura 2020 T2A7
Gegeben sind die Punkte $A(1,0,1), B(7,-9,4)$ und $C(3,0,-2)$.\newline Bestimme die Punkte auf der Gerade durch $A$ und $B$, die vom Punkt $C$ den Abstand $5$ haben.
\vspace{5pt}
\antwort{21.5}

%---------------------------------------------------------
\newpage
\noaddpoints
\question[\totalpoints] %Matura 2022 A14
\addpoints
In der Grundebene (xy-Ebene) befinden sich drei Dörfer bei den Koordinaten $A(9,2,0), B(1,8,0)$ und $C(9,12,0)$. Zwei Berggipfel haben die Koordinaten $D(5,4,6)$ und $E(3,6,8)$.\newline
\begin{minipage}{0.6\linewidth}
    \begin{parts}
        \part[2] Zeige durch Rechnung, dass das Dorf $A$ und die beiden Berggipfel $D$ und $E$ nicht auch einer Gerade liegen.
        \part[3] In welchem Punkt $P$ der Grundebene $ABC$ befinden sich beide Bergspitzen $D$ und $E$ genau hintereinander?
    \end{parts}
\end{minipage}
\hfill
\begin{minipage}{0.38\linewidth}
    \centering\includegraphics[scale=0.8]{Vektorgeometrie/data-vektor/dörferberge.png}
\end{minipage}
\vspace{5pt}
\antwort{17.5}

\end{questions}
\end{document}